\section{Inverse-Lucas correction: 85 operations in ten positive variables}
\label{sec:inverse-lucas}

\paragraph{Result and scope.}
The literal polynomial \texttt{positive/equation\_bV.json} has ten
independently strictly positive integer unknowns, exact degree \(304\),
and \(85\) live binary arithmetic operations.  The newest inspected
Lucas-correction parent has the same degree and \(87\) operations.
The square-weight control has \(84\) operations and degree \(410\);
it is dominated by known earlier points and is not promoted as a frontier
improvement.  The complete indexed BFS943/radix67 tuple is unchanged,
with largest fixed port of \(592{,}884\) bits.

Composed with the explicitly retained ordinary native/refutation
calibration, the proof below gives
\[
 \mathrm{ZFC}\vdash
 \bigl(\neg\exists\boldsymbol{x}\in\mathbb Z_{>0}^{10}\;
 P(\boldsymbol{x})=0\bigr)\longleftrightarrow\operatorname{Con}(\mathrm{ZFC}).
\]
This neither proves consistency nor asserts a zero.  It is an ordinary
mathematical proof with exact finite artifact checks, not a new end-to-end
proof-assistant object or a global-minimum claim.

The new mechanism replaces the forward rank correction by an inverse
Lucas correction at the existing abscissa \(k\).  This saves two
multiplications and removes the exceptional negative-rank transport cases.
A genuine obstruction at the prime \(3\) means that the old denominator
proof cannot simply be reused.  Projection must precede full denominator
recovery; a new absolute bound handles a possibly negative corrected quotient
in the positive lift.  The second-order expansion and the projection-first
strategy themselves are explicitly inherited from the inspected Lucas
Corrections and Cofinal Witnesses archives at commit
\texttt{7a8da414376a80b90ecf799f37ceb26fb9ac674e}.

\subsection{Complete polynomial}
The only unknowns are
\[
 J,F,\alpha,i,T,w,s,n,\eta,\zeta\in\mathbb Z_{>0}.
\]
The fixed integers \(\beta,K_0,M_C,M_F\) are the whole supplied tuple,
with \(\beta\ge31\) odd, \(K_0>0\), and
\(0<M_C<\beta<M_F<2\beta\).  Define polynomial abbreviations:
\begin{gather*}
 h=\beta J,\quad q=h+1,\quad C=q-F-\alpha,\\
 X=qw,\quad Y=qs,\quad k=\eta+\zeta,\quad c=kY+\eta,\\
 A=Y(X+1)+1,\quad \Delta=A^2-1,\quad H=4A-5,\\
 v=XY^2,\quad \ell=vk=(XY)(kY),\quad E=1+\ell(\ell+k),\\
 R=[h(q-F)+\alpha](q^2-1)+(M_C+qM_F)J,\\
 B=(K_0+w)C-F,\quad U=ic^2,\quad a=\Delta U^2E,\\
 F_0=1+a,\quad G=EF_0,\quad F_s=1+4aF_0,\quad Q=\Delta(F_s-1),\\
 b=cT-R,\quad V=F_sb-c,\quad K=X(3c-2X)+2,\\
 e_6=6(k-R-1)-\ell(k^2-4),\quad D_6=6q^6,\quad M=hD_6H.
\end{gather*}
Use \(L=bV\) for the principal equation, or \(L=V^2\) for the control,
and put
\begin{align}
 Z&=hD_6\bigl[H(n-LB)+K\bigr]
       -H\bigl[D_6B+h(e_6-6q^4s)\bigr],\label{eq:inverse-Z}\\
 P&=Q\bigl(Z^2-M^2GV^2\bigr)+(M^2G-Z^2).\label{eq:inverse-P}
\end{align}
The single equation is \(P=0\).  All displayed powers have fixed
nonnegative integer exponents.  The DAG does not compute \(v\)
separately: it reuses \(\ell=(XY)(kY)\).  There are no division gates,
variable exponents, sequence functions, or additional quantified variables.

For proof only define
\begin{gather*}
 e_*=\frac{e_6}{6}=k-R-1-\frac{vk(k^2-4)}{6},\quad
 \rho=\frac{e_*}{q^6},\quad\sigma=\frac{s}{q^2},\\
 t=\frac{B}{h},\quad W=1+hL,\quad
 \kappa=\frac{K}{H},\quad z=\frac{Z}{M}.
\end{gather*}
Without assuming any integrality of these quotients,
\begin{equation}
 z=n+\kappa+\sigma-Wt-\rho.\label{eq:inverse-carrier}
\end{equation}

\subsection{Root extraction and uniform signed transport}
The unchanged root and folding arguments are retained in
\texttt{inherited\_lucas/PROOF.md}, Sections 2--3, with their ordinary
Pell sources.  We state their premises and conclusions to specify the
interface precisely.  Positive inputs give \(q\ge32\), \(X,Y\ge q\),
\(k\ge2\), \(c>k\), \(A>4\), \(U\ge c^2\), \(F_s>2c\),
and \(Q>c^2\).  Moreover
\[
 F_s=(1+2a)^2,\qquad Q=4\Delta^2U^2G,\qquad\gcd(E,1+a)=1.
\]
At a zero, prime valuations in
\((Q-1)z^2=G(QV^2-1)\) give \(G=g^2\), \(E=\tau^2\),
\(F_0=f_0^2\), with \(g=\tau f_0>0\) integral.  Rational adjacent-Pell
descent makes \(y=|z|/g\) integral, and
\[
 QV^2-(Q-1)y^2=1,\quad z=\varepsilon gy,\quad
 \varepsilon\in\{1,-1\}.
\]
Writing \(f=1+2a\), \(I=2Ug\), we have
\(f^2-\Delta I^2=1\), \(Q=(\Delta I)^2\), and \(c^2\mid I\).
Signed all-index folding gives
\[
 c=\psi_A(p)\quad(p>0\text{ odd}),\qquad
 k=2\psi_{1+2v}(m),\quad m\ge2,
\]
where \(\psi_a(0)=0\), \(\psi_a(1)=1\), and
\(\psi_a(j+2)=2a\psi_a(j+1)-\psi_a(j)\).
In particular \(k\ge8v+4\), \(c>8XY^3\ge8q^4\), \(c>X\), and
\(K>0\).

The inherited exclusions of main and auxiliary index three use only
\(q\mid X,Y\), not the rank correction.  They give, for either weight,
\begin{equation}
 b\ne0,\quad \operatorname{sign}V=\operatorname{sign}b,
 \quad |V|>|b|,\quad L>0,\quad gy<L/2,
 \quad {|b|\over q^6}<L/2.\label{eq:inverse-abs}
\end{equation}
For \(L=bV\), the main index-three remainder would force
\(2Y-X\in\{1,2\}\), impossible modulo \(q\).  The two signed auxiliary
index-three possibilities would give
\[
 F_s(b-4\Delta)=c-(4\Delta+3),\qquad
 F_s(b+4\Delta)=c+(4\Delta+3),
\]
whose right sides lie strictly between \(0\) and \(F_s\).
Thus the auxiliary index is at least five, giving \(|b|>4g\) and
\(gy<3L/8\).  Neither sign branch is discarded.

The new correction has an advantage before any sign of \(R\) is known:
\(k\ge2\) makes \(vk(k^2-4)/6\ge0\).  Therefore
\[
 e_*\le k-R-1<b,\qquad \rho<|b|/q^6.
\]
If \(B\le-1\), then \(Wt<-L\), and
\eqref{eq:inverse-carrier}--\eqref{eq:inverse-abs} force
\[
 n=\varepsilon gy-\kappa-\sigma+Wt+\rho
   <L/2-L+|b|/q^6<0,
\]
a contradiction.  This handles all negative \(R\) and \(R=0\)
uniformly.  Hence \(B\ge0\), so \(C>0\).  Only now is the inherited
mask bound \(3q<R<q^4-2\) invoked.  Then \(c>2R\) and \(b,V>0\);
the normalized-rank and signed auxiliary-index theorems give \(p=R\).
Thus \(R\) is positive odd; mask parity gives even \(q\) and odd \(J\).

\subsection{Inverse-Lucas authentication at all primes}
The inherited, universally valid finite Lucas expansion gives
\begin{equation}
 k=2m+8vC_m+32v^2Z_m,
 \quad C_m=\binom{m+1}{3},\quad Z_m\in\mathbb Z_{\ge0}.
 \label{eq:inverse-expansion}
\end{equation}
Binomial notation occurs only in the proof.  Since \(k=2u\) is even,
\(k(k^2-4)/6=4u(u-1)(u+1)/3\) is integral.  Thus \(e_*\) is
integral at a zero.  Clearing the integral carrier gives
\(6q^6\mid hH(e_6-6q^4s)\); cancelling the ordinary factor six and
using \(h=q-1\), \(H\equiv-1\pmod q\), gives
\begin{equation}
 q^6\mid e_*-q^4s,\qquad q^4\mid e_*.
 \label{eq:inverse-q4}
\end{equation}
No modular inverse of six is used.

Put \(t_0=2m\), \(d=2m-R-1\), \(A_m=C_m+4vZ_m\), and
\(S_m=k^2+kt_0+t_0^2-4\).  The exact inverse identity is
\begin{equation}
 3(e_*-d)=4v^2(24Z_m-A_mS_m).
 \label{eq:inverse-identity}
\end{equation}
Indeed \(t_0(t_0^2-4)=48C_m\), \(k-t_0=8vA_m\), and
\(k(k^2-4)-t_0(t_0^2-4)=(k-t_0)S_m\).
Since \(v=q^3ws^2\), the right side is divisible by \(q^6\).
For each prime \(\lambda\mid q\), with \(a_\lambda\ge1\) its
valuation in \(q\),
\[
 \operatorname{ord}_\lambda(e_*-d)
 \ge6a_\lambda-\operatorname{ord}_\lambda(3)\ge4a_\lambda.
\]
Hence \(q^4\mid d\).  The first base \(1+2XY^2\) exceeds \(A\).
If \(m\ge R\), monotonicity would give
\(k/2\ge\psi_A(R)=c\), contrary to \(c>k\).  Thus \(1\le m<R\),
and \(|d|<R<q^4\), proving
\begin{equation}
 d=0,\qquad m=(R+1)/2.\label{eq:inverse-index}
\end{equation}

It would be false to infer full \(q^6\)-divisibility here.
For \(q=3,v=27,m=2,R=3\), the inverse error is
\(-47{,}911{,}824\equiv243\pmod{729}\).
This modular diagnostic is not a purported zero of the full polynomial.
Evenness of \(q\) alone does not remove the obstruction at three.

\subsection{Projection first, denominator second}
Because \(z=Z/M\) is integral, reducing the literal numerator
\eqref{eq:inverse-Z} modulo \(H\) gives \(H\mid6hq^6K\).
As \(\gcd(q,H)=1\),
\begin{equation}
 H\mid6hK.\label{eq:inverse-project}
\end{equation}
This uses neither integral \(\rho\) nor integral \(\sigma\), and does
not cancel three.  Write \(R=2r+1\), so \(r\ge48\).
The retained uniform ratio theorem, licensed by the authenticated
indices, gives
\[
 Y>X^r/3,\qquad X^{r+1}<3Y(X+1)<H.
\]
Modulo odd \(H\), the scaled recurrence gives
\(3\cdot2^{R-1}\psi_A(R)\equiv4^R-1\), because \(4A\equiv5\).
Consequently \(2^{R-1}K\equiv-(X-2^R)(2^RX+1)\pmod H\).
Equation \eqref{eq:inverse-project} therefore implies divisibility
of this whole product by \(H\), with multiplier \(6h\).

For an explicit height check let \(u_0=2^R=2\cdot4^r\).
Since \(6<q^3\) and \(q^3h<X^4\),
\begin{align*}
 6h\,|(X-u_0)(u_0X+1)|
 &<X^4(X+u_0)(u_0X+1)\\
 &\le4\cdot4^rX^6+8\cdot16^rX^5<X^{r+1}<H.
\end{align*}
After division by \(X^{r+1}\), the two terms decrease with
\(X\ge32,r\ge48\); their sum at \((32,48)\) is
\(2^{-117}+2^{-25}<1\).  Hence the whole divisible product is zero,
and its other factors are positive.  Thus \(X=2^R\); the same universal
projection then gives integral positive \(\kappa=K/H\).

Now \(q\mid X\) implies \(q=2^a\), \(a\ge5\).
Equation \eqref{eq:inverse-identity}, with \(d=0\), finally gives
\(q^6\mid e_*\), because three is a unit modulo \(q\).
Thus \(\rho\) is integral.  The carrier says
\(\sigma-WB/h\) is integral.  Multiplying by \(q^2h\) and reducing
modulo \(q^2\) gives \(q^2\mid s\).  Hence \(q^3\mid Y\),
\(\sigma\) is integral, and \(h\mid WB\).  Since \(W\equiv1\pmod h\),
\(h\mid B\), so \(t\) is a nonnegative integer.

The fixed-minus parity, exact half-binomial rounding, and fresh original
native reconstruction now have their full retained premises.
That source uses base \(Y(X+1)+2\), not the auxiliary base \(A\).
The unchanged whole-tuple calibration produces a finite ZFC refutation.

\subsection{Constructive positive completeness}
A finite refutation supplies the retained canonical outer data
\(q=2^a\), \(X=2^R=qw\), \(R>3q\), \(R\equiv3\pmod4\),
\(q^3\mid Y\), \(C>0\), and \(t=B/h\ge0\) integral.
Set \(s=Y/q>0\), and form
\[
 c=\psi_A(R),\quad k=2\psi_{1+2XY^2}((R+1)/2),\quad
 \tau=\chi_{1+2XY^2}((R+1)/2),\quad
 \eta=c-kY,\quad\zeta=k-\eta.
\]
The exact ratio theorem gives \(0<\eta<k\), hence both last inputs
are positive.  Equation \eqref{eq:inverse-identity} and dyadic \(q\)
give integral \(\rho\).  Projection gives positive integral \(\kappa\),
and \(\sigma=Y/q^3\) is integral.

A positive period of the Pell matrix modulo \(c^2\tau\) gives a seed
\((f_0,I_0)\) with \(c^2\tau\mid I_0\).
Set \(i=I_0/(c^2\tau)>0\).  Then \(a=\Delta I_0^2=f_0^2-1\),
\(F_0=f_0^2\), and \(G=(\tau f_0)^2\).
The doubled pair \(f=2f_0^2-1\), \(I=2f_0I_0\) has exactly the
required definitions.  The retained auxiliary-completeness theorem gives
positive \(V,y\) and the positive integer
\(T=(V+c+Rf^2)/(cf^2)\), the actual unshifted input.

The new \(\rho\) may be negative; the old nonnegative-quotient argument
is not reused.  Since \(k\ge R+1\) and \(vk^2>6\),
\[
 |\rho|\le {k-R-1+vk(k^2-4)/6\over q^6}
       <{k+vk^3/6\over q^6}<vk^3/3.
\]
For every permitted \(i\ge1\),
\[
 g>\sqrt\Delta\,c^2E>c^2v^2k^2>vk^3,
\]
so \(|\rho|<g/3\).  Also \(\sqrt\Delta>A-1=Y(X+1)>X\) and
\(cE>4\), hence \(g>4Xc\).  As \(\kappa<Xc\) and
\(\sigma\le Y<Xc\), we have \(\kappa+\sigma<g/2\).
For either positive weight put
\[
 n=gy+(1+hL)t+\rho-\kappa-\sigma.
\]
This is integral, and \(y\ge1\), \(t\ge0\) imply
\(n>g-g/3-g/2>0\).  Its carrier is \(z=gy\), so \(P=0\).
All ten independent coordinates have been constructed strictly positively;
no inverse positive shear or unproved sign is used.  Every permitted
period seed works, without additional cofinal enlargement.

\subsection{Exact costs and verification}
Deleting the incoming gate \(v=(XY)Y\), restoring the sharing
\(\ell=(XY)(kY)\), and replacing its four-gate rank cubic by
\(k^2\), \(k^2-4\), \(\ell(k^2-4)\) saves two multiplications.
The principal has \(51\) multiplications, \(18\) additions and
\(16\) subtractions.  All inputs, ports and rows are live.

Put \(q_0=\beta J\), \(k_0=\eta+\zeta\),
\(C_0=q_0-F-\alpha\), and
\(B_0=k_0sT-q_0^2(q_0-F)\).  The \emph{complete} highest homogeneous
forms are
\begin{align*}
 [P_{bV}]_{304}
 &=-9\cdot2^{12}q_0^{128}w^{34}s^{78}i^{14}k_0^{44}B_0^2,\\
 [P_{V^2}]_{410}
 &=9\cdot2^{16}q_0^{166}w^{46}s^{104}i^{20}k_0^{60}B_0^4C_0^2.
\end{align*}
Their full expansions have respectively \(274\) and \(3750\) terms.
Nonzero coefficients at maximal \(T\)-power establish the exact degrees
for the actual \(\beta>0\), not merely a formal upper bound or a ray sample.

The producer-independent checker covers all actual rows with \(37\)
exact free-polynomial cuts per graph, verifies full leading forms and
every coefficient byte, reconstructs the four fixed numerals, and rejects
both alternative operations at every gate plus a missing-correction
mutation: \(171\) rejections for the principal and \(169\) for the control.
Separate finite recurrence controls include the obstruction at three;
they do not replace the unbounded arithmetic proof.
Normal and optimized Python receipts agree byte for byte.
\begin{verbatim}
python3 -B positive/build.py --check
python3 -B positive/verify.py
python3 -O -B positive/verify.py
\end{verbatim}

\subsection{Whole complemented-label radix47 sidecar}
The same two free-port row lists are separately specialized to the
\emph{entire} frozen complemented-label radix47 tuple.  The principal
remains \(85/304\), but its largest fixed port is \(415{,}904\) bits and
the four raw numerals occupy \(207{,}214\) bytes rather than \(295{,}388\).
This reuses the incoming encoding and controller; it is not a new radix
theorem and is not the repository's conditional three-anchor radix47 layout.
The canonical radix67 equations remain separate.

The new arithmetic is uniform in the same scalar, source-shift mask and
population interface already checked for the whole47 tuple.  In soundness
it recovers \(C>0\), \(X=qw=2^R\), \(q^3\mid Y\), the native
half-binomial, and \((K_0+w)C-F=t(q-1)\), \(t\ge0\).
These are exactly the retained whole47 bridge's premises, with
\(W_{\rm native}=w\) and source transport variable \(t+1\).
Its fresh Pell witnesses are constructed at \(Y(X+1)+2\), before the
actual complemented-label decoder is applied.  Conversely, the actual
whole47 controller supplies dyadic \(q\) and the same outer interface;
choosing \(w=W_{\rm native}\), \(s=Y/q\), and applying the explicit
positive construction above gives a new zero.  No old67 decoder or
controller is silently substituted.

The complete ordinary native bridge, incoming P66 equation, all four
payloads, \(3297\) coefficients, \(51\) native target identities, and its
\(37\) pinned dependencies are retained under
\path{positive/radix47/inherited/}.  The sidecar checker requires identical
free-port rows and an exact whole-tuple binding, reconstructs every numeral,
and repeats the cuts, leading-form and mutation checks.  Its normal and
optimized receipts agree.  The default main checker includes this sidecar;
its literal binding can also be reproduced independently:
\begin{verbatim}
python3 -B positive/radix47/build.py --check
python3 -B positive/radix47/verify.py
\end{verbatim}
